\section{Five Layers of Engineering Semantics}
\label{sec:layers}
The layers organise sources by their role in agent grounding, rather than
replacing lifecycle or hierarchy axes in industrial reference architectures.
One standard can contribute several capabilities, and a deployment can combine
them without materialising all five layers in one graph. Specifications and
reference documentation provide the survey evidence; exposed implementation
profiles must be assessed separately.

\subsection{L0: Identification and property dictionaries}
ECLASS, IEC common data dictionaries and industrial reference identifiers link
an engineering property to a stable definition; GS1 identifiers distinguish
products and individual assets across systems. This is useful when two suppliers
use different labels for the same quantity, or the same label for different
quantities. An agent should resolve to an identifier whose mapping is inspectable,
not simply to the closest English description. Concept identity and instance
identity are distinct: identifying ``pressure'' does not identify the pressure
transmitter on a particular filler. The latter needs asset membership and a
specific observation source. Persistent identifiers also require governed mappings
when hardware is replaced; stability of a string alone does not establish
continuity of the physical asset \cite{eclass,iec61360,iec61987,iso29002}.

The practical mechanism is a versioned mapping from source identifiers to asset
instances and property concepts. A mapping should identify the responsible source,
validity interval and replacement history. Dictionary lookup reduces ambiguity
between vendors, but cannot resolve an observation's current owner without an
instance model. Hence L0 supports reference integrity while L1 and observation
metadata complete the operational referent.

\subsection{L1: Information and interface models}
OPC UA combines typed address spaces, attributes, access rights and callable
methods; companion models add domain content \cite{iec62541}. AAS supplies
submodels, semantic identifiers and structured asset descriptions
\cite{idta2023aas}. AutomationML supports engineering exchange, ISA-95 describes
equipment hierarchy, and MTP and PackML contribute services and behaviour.
WoT Thing Descriptions and NGSI-LD provide additional interface and entity
representations. These technologies make introspection and tool construction
possible \cite{iec62714,iec62264,vdi2658,isa88,w3c2023wot,etsi2023ngsild,iec61850},
but the content actually exposed must be audited. An agent cannot
check an unpublished interlock simply because the host standard is extensible.

Conversely, it is inaccurate to describe all information models as static or
incapable of behaviour: state-machine and method information can be represented
in this layer. The bottleneck is often an engineering publication decision.
Controllers, alarm configurations and recipes already contain constraints, but
their relation to a supervisory command is not necessarily delivered through
the agent-facing interface. Typed access is a necessary integration baseline,
not sufficient evidence of a complete operating envelope.

The integration must also distinguish access rights from operating permission.
A point can be technically writable yet unavailable in the current state or
outside an operator's scope. An engineering range is not necessarily a recipe
limit or a safety limit. Collapsing these into a single numeric boundary loses
the reason for a rejection and makes changes difficult to audit. Model versions
and the origin of each constraint therefore belong in the agent evidence contract.

\subsection{L2: Formal semantics, validation and query}
RDF provides globally named relations and a common graph substrate. OWL can
support class and relation entailment, while SHACL checks a finite data graph
against explicit shapes \cite{knublauch2017shacl}. The distinction matters:
open-world logical consistency does not by itself enforce a closed set of
mandatory action parameters or establish that all interlocks have been supplied.
A deployment must declare which constraints are enforced, what missing data
means and which reasoning profile is supported.

SPARQL executes joins, traversals, filtering and aggregation outside the language
model, returning a bounded result rather than every supporting record. PROV-O
can associate transformations and observations with sources and versions. Neither
a successful query nor shape conformance establishes the correctness of the
engineering mapping. Graph construction, source calibration, shape completeness
and unit fidelity remain upstream obligations. Declarative constraints can be
shared and inspected, but a competent procedural implementation may enforce
exactly the same rules; the R1 control makes that comparison explicit.

Validation concerns the proposed action and its referenced snapshot, not only
data quality of the underlying graph. Lifting a proposal into a small action
graph makes target, value, unit, scope and evidence time explicit. A report can
then distinguish signature errors, incompatible quantities, violated guards and
missing observations. A repair loop should preserve the original authorization
and provide a bounded opportunity to correct the proposal, rather than weakening
the constraint after repeated refusal.

\subsection{L3: Quantities, observations and domain relations}
QUDT supplies quantity kinds, dimensions and conversion metadata \cite{qudt}.
The distinction between a display symbol and a conversion definition is crucial:
``bar'' and ``psi'' name compatible units but do not themselves state the
conversion coefficient. SOSA/SSN distinguishes a sensor, an observed property,
an observation and its result \cite{haller2019ssn}. This permits an answer to
cite the measurement and acquisition time it claims to report, rather than
present a number detached from a source.

Brick supplies typed equipment relations \cite{balaji2016brick}; SAREF and its
extensions describe devices and domain functions; DEXPI carries process-engineering
structure. These should be selected for actual industrial fit rather than treated
as interchangeable taxonomies. Flow connectivity, containment and observation
provenance are different relations. A compressor feeding several units supports
impact analysis, but a flow graph alone does not establish the causal mechanism
of a recorded failure. Root-cause analysis needs incident evidence and a validated
causal/diagnostic model beyond generic graph traversal.

Observation quality and freshness must also be preserved through conversion.
Expressing a stale sample in another unit does not make it current, and retaining
a value without its source quality loses the reason an operator might decline
to act. Engineering tolerances belong to the comparison procedure: matching
observations exactly is inappropriate when sensor resolution, rounding and
calibration uncertainty are material. The reference benchmark's numeric tolerance
is therefore a reproducible experimental setting, not a universal industrial rule.

\subsection{L4: Cross-organisational policy and trust}
Dataspaces add participant identity, catalogue metadata, exchange agreements and
usage-policy mechanisms \cite{idsa2022ram}. Catena-X/Manufacturing-X highlight
versioned shared semantic models; ODRL expresses permissions, prohibitions and
duties \cite{w3c2018odrl}. A supplier observation may be structurally interpretable
yet unsuitable for a particular purpose or untrusted. Policies are therefore a
separate gate from physical operating constraints. Policy expressions do not
automatically enforce downstream confidentiality, and keeping information in an
agent's memory is not universally equivalent to redistribution. Enforcement
depends on the policy, processing boundary and trusted runtime. L4 is included
for architectural completeness and is not experimentally validated here.